\begin{examplebox}
\textbf{例1：先标准化。} $9x^2-16y^2=144$ 化为 $x^2/16-y^2/9=1$，
所以 $a=4,b=3,c=\sqrt{16+9}=5$，焦点为 $\boxed{(\pm5,0)}$。
回验：右顶点 $(4,0)$ 到两焦点距离为 $9,1$，差为 $8=2a$。

\medskip
\textbf{例2：逆向求方程，要补齐位置条件。}
已知\textbf{中心在原点、焦点在坐标轴上}的双曲线有焦点 $(5,0)$，且实半轴长 $3$。
于是 $c=5,a=3$，$b^2=25-9=16$，正项为 $x^2$，得
\[
 \boxed{\frac{x^2}{9}-\frac{y^2}{16}=1}.
\]
回验：焦点 $c^2=9+16=25$，实轴长为 $6$。
只给“一个焦点是 $(5,0)$”而不限定中心与方向，不能唯一确定标准方程。

\medskip
\textbf{例3：同焦点不等于同参数。}
双曲线与椭圆 $x^2/25+y^2/9=1$ 同焦点，实轴长为 $6$，求方程。
\begin{center}\begin{tikzpicture}[font=\small,>=stealth]
\draw[black!30,rounded corners=2pt] (0,0) rectangle (14.4,8.2);
\node[anchor=west,font=\small\bfseries] at (.5,7.72) {同一对焦点，椭圆与双曲线的 $a,b$ 仍然不同};

\begin{scope}[shift={(7.2,4.5)},scale=.48]
\draw[black!35,->] (-6,0)--(6,0) node[right] {$x$};
\draw[blue,very thick] (0,0) ellipse (5 and 3);
\draw[orange!85!black,very thick,dashed,domain=-4.5:4.5,samples=121,variable=\t] plot ({3*sqrt(1+\t*\t/7)},\t);
\draw[orange!85!black,very thick,dashed,domain=-4.5:4.5,samples=121,variable=\t] plot ({-3*sqrt(1+\t*\t/7)},\t);
\fill (-4,0) circle (2pt) node[below=5pt] {$F_1$};
\fill (4,0) circle (2pt) node[below=5pt] {$F_2$};
\end{scope}
\node[blue] at (7.2,1.65) {椭圆：$a=5,b=3$；$c^2=25-9=16$};
\node[orange!85!black] at (7.2,1.1) {双曲线：$a=3,b=\sqrt7$；$c^2=9+7=16$};

\node[anchor=west] at (.7,.55) {同尺度、共用焦点 $(\pm4,0)$；同焦点只固定 $c$，并不固定 $a,b$};
\end{tikzpicture}\end{center}

椭圆焦点为 $(\pm4,0)$，故双曲线 $c=4$；$2a=6$ 给出 $a=3$，
于是 $b^2=16-9=7$，答案为 $\boxed{x^2/9-y^2/7=1}$。
回验：$9+7=16=c^2$，右顶点到两焦点距离为 $7,1$，差 $6$。
若误取实半轴 $a\ge4$，就不能得到 $b>0$ 的这类双曲线。
\end{examplebox}
\clearpage
